\subsection{An arithmetic pro-two replacement and simultaneous relative-prime framing}

\begin{statement}{D.46}{Proved}{A replacement with the prescribed extra CM pair on the full relative pronilpotent source}
Retain \(N=\ker(B\to Q)\), with \(Q\) of order \(576\)
and \(N\) free of rank \(97\). For every prime \(p\), define
the actual source kernels and relative completions
\[
 R_p=\ker(\widehat N\to N^{(p)}),\qquad
 \mathcal P_p=\widehat B/R_p,\qquad
 R_{\rm nil}=\bigcap_pR_p,\qquad
 \mathcal U=\widehat B/R_{\rm nil}.
                                                               \tag{D.46.1}
\]
There are exact marked identifications
\[
 \mathcal U=\prod_Q\mathcal P_p,\qquad
 \ker(\mathcal U\to Q)=\prod_pN^{(p)}.                          \tag{D.46.2}
\]
In particular this retains the entire nonabelian relative
pro-\(2\) kernel, not a finite-degree or class-two quotient.

There is a continuous peripheral automorphism \(\alpha^\dagger\)
of \(\mathcal U\), with multiplier one and an actual
\(\widehat F'\) parameter, such that its component at every
odd prime is the identity, while its component at two
is induced by one actual Galois commutator \(h\).
Its finite \(G_6=B/[N,N]N^6\) shadow has
\[
                 V_{12}=I,\qquad T_4=I,\qquad W_3=[i],
                                                               \tag{D.46.3}
\]
with the exact global-sign and conjugator conventions of D.25.
Thus it realizes the prescribed extra CM pair on this whole
relative-pronilpotent quotient. Its finite refinements retaining
\(G_6\) give a compatible family of nonarithmetical gentle
GT shadows. No assertion of full profinite GT0 realization
is made by identifying this quotient family.

The two-component action is particularly explicit.
Put \(\zeta=\zeta_3\) and
\[
 a=1-2\zeta^2,\qquad
 \eta_*=\frac a{\bar a}=\frac{-3-8\zeta^2}{7}.
                                                               \tag{D.46.4}
\]
The normalized action of \(h\) on the full sphere group
\(F^{(2)}\) is the identity. Its actual action on the compact
cubic elliptic Tate module is multiplication by \(\eta_*\).
In the elliptic graph basis of D.45 its matrix modulo
sixteen is
\[
                  \begin{pmatrix}11&8\\8&3\end{pmatrix}.       \tag{D.46.5}
\]
The arithmetic element \(h\) itself has the CM pair
\((-I,[i])\); (D.46.3) is recovered from the new joint
relative-prime action through \(G_2\) and \(G_3\).
The proof specifies every conjugator map producing that
change of framing, while keeping the two-component
automorphism unchanged.

There is a canonical identification
\(\mathcal P_2\simeq\mathcal P_2^{\,K}\), where
\(\mathcal P_2^{\,K}\) is the full rank-four relative
source of D.34. The projection of this arithmetic
component to (D.45.3) is
\[
             (1_{\mathcal P_F},\alpha_{h,E}).                  \tag{D.46.6}
\]
Its lift through the rank-four mixed-commutator kernel
is the actual restriction of \(h\). Thus this particular
replacement solves that entire relative pro-\(2\)
lifting step. It replaces the excluded D.33/D.39
lower action: the old sphere action is already
nonidentity in degree seven, while (D.46.6) is
the identity on the whole sphere factor.

The parameter is the exact compatible tuple
\[
                 f^\dagger=(f_{h,2},1,1,\ldots)
                    \in\operatorname{im}(\widehat F'
                                            \to\mathcal U).
                                                               \tag{D.46.7}
\]
It has trivial image in every pure pro-\(p\)
completion of \(F\), but it is nontrivial in
\(\mathcal U\), as witnessed by (D.46.4).
The exact projection kernels are
\[
 \ker(\mathcal U\to\mathcal P_p)=\prod_{q\ne p}N^{(q)},
 \qquad
 \ker(\widehat B\to\mathcal U)=R_{\rm nil}.                     \tag{D.46.8}
\]
Under the actual marked epimorphism to \(S_7\)
whose permutations, on \(0,\ldots,6\), are
\[
 s\mapsto(1,0,3,2,5,4,6),\qquad
 r\mapsto(3,6,2,5,1,0,4),
                                                               \tag{D.46.9}
\]
the latter source kernel maps exactly onto \(A_7\).
The relative centralizer elements used below do not
come from the corresponding full \(\widehat B\)
centralizers. These two exact facts identify the
remaining full-source lifting problem.

This continues Part D(e), section 51.16, folio
734 verso, scan 174, and the two-torus follow-up
to LM2, section 51.15, folios 725--727,
scans 156--160.
\end{statement}
\begin{proof}
\emph{The actual arithmetic commutator.}
Write \(k=\mathbb Q(\zeta)\) and
\(\mathcal O=\mathbb Z_2[\zeta]\).
D.35 proves that the actual image of
\(G_k\) on \(T_2E_3\), for
\(E_3:y^2=x^3+1\), is
\(U=1+2\mathcal O\). Its full division
field over \(k\) is therefore pro-two.
The cyclic cubic field
\(k(\sqrt[3]{4})\) is disjoint from it.
Consequently the joint map
\[
 G_k\longrightarrow U\times C_3,\qquad
          g\longmapsto(\rho_{E_3,2}(g),\kappa_{4,3}(g))
\]
is onto. Choose an actual \(g\) with
image \((a,0)\), where \(a\) is
(D.46.4). This is a use of the proved
arithmetic image and disjointness,
not an assumed realization of a
matrix. The CM norm is \(a\bar a=7\),
so \(\chi_2(g)=7\). Also
\(\chi_3(g)\equiv1\pmod3\) and
\(\kappa_{2,3}(g)\equiv0\pmod3\).

Since \(k/\mathbb Q\) is ramified
at three, choose \(\tau\) in a
fixed inertia group \(I_3\) that
induces the nontrivial automorphism
of \(k\). Set
\[
                   h=[g,\tau]=g\tau g^{-1}\tau^{-1}.           \tag{D.46.10}
\]
The entire cyclotomic character
of \(h\) is one, because the
character is abelian. The
normalized pro-two action on
the thrice-punctured line is
unramified at three. The exact
statement for the tangential
automorphism action appears in
\cite[\S8.5, preceding Theorem 8.7]{E_HainMatsumoto},
with \(N=1,\ell=2\); \S8.4
states the corresponding outer
action for every prime.
Thus \(\tau\), and hence the
commutator \(h\), acts trivially
on \(F^{(2)}\).
The same conclusion follows
from the outer statement and
the normalization argument of
D.23: an inner action fixing
\(x\) has conjugator in
\(\langle x\rangle\), and
comparison with
\(y\mapsto f^{-1}yf\),
\(f\in(F^{(2)})'\), forces
both abelian exponents to
vanish.

The exact affine Kummer
commutator identity is
\[
 \kappa_b([g,\tau])
    =(1-\chi(\tau))\kappa_b(g)
                       +(\chi(g)-1)\kappa_b(\tau).             \tag{D.46.11}
\]
Modulo three, the two
conditions imposed on \(g\)
make \(\kappa_{2,3}(h)=0\).
At two, \(\chi_2(\tau)=1\).
Also \(\kappa_{2,2}(\tau)=0\):
every \(2\)-power root of the
\(3\)-adic unit \(2\) lies
in the maximal unramified
extension, by Hensel's lemma
and its prime-to-three
derivative. Therefore
\(\kappa_{2,2}(h)=0\).
In particular D.18's actual
matrix gives
\[
 V_h=\begin{pmatrix}1&-8\kappa_2(h)\\0&1\end{pmatrix},
                         \qquad V_h\bmod12=I.                 \tag{D.46.12}
\]
The induced normalized action
on \(Q\) is the identity.

The compact cubic cover
\(\mathcal E_0\) of D.34 has
the exact twist
\(\rho_{\mathcal E_0,2}
=\zeta^{-\kappa_{4,3}}\rho_{E_3,2}\).
Its value on \(g\) is \(a\),
and conjugation by the
action of \(\tau\) conjugates
the CM coefficient ring.
Hence its value on \(h\)
is exactly \(a/\bar a=\eta_*\),
regardless of the particular
semilinear matrix of \(\tau\).
Equation (D.46.11) also gives
\(\kappa_{4,3}(h)=0\), so the
actual \(E_3\) scalar is the
same. Its reduction modulo
four is \(-I\).
The Galois comparison of
D.25--D.36 then gives
the actual arithmetic pair
\[
                         (T_{h,4},W_{h,3})=(-I,[i]).          \tag{D.46.13}
\]
This use of the cover-to-fixed-fibre
map is confined to the actual
Galois element \(h\).

\emph{The full rank-four and two rank-two actions.}
The subgroup \(N\) has
\(2\)-power index \(32\)
in \(K\). The full induced
pro-two topology argument
of D.45 proves equality
of the actual source kernels
\[
 \ker(\widehat N\to N^{(2)})
                =\ker(\widehat K\to K^{(2)}).                 \tag{D.46.14}
\]
Thus the relative completions
over \(Q\) and \(Q_0\) are
canonically the same marked
group. The Galois action
preserves \(N\) because it
acts on \(Q\) by (D.46.12),
and preserves its pro-two
residual. It consequently
induces an actual automorphism
\(\alpha_{h,2}\) of this
entire relative source.
Restriction and the two
actual quotient maps of
D.45 give (D.46.6).
The original arithmetic
automorphism is its lift
through the exact mixed
commutator kernel, so the
lifting assertion does not
follow merely from the
homology determinant.

In the elliptic free basis
\(e_1,e_2\) of D.45, the
literal graph conjugation
by \(t^2\) has matrix
\[
                      Z_E=\begin{pmatrix}0&1\\-1&-1\end{pmatrix}.
                                                               \tag{D.46.15}
\]
It is the image of
\(\zeta^2\) under the
same CM identification as
in D.35. Substitution into
\((-3I-8Z_E)/7\) gives
(D.46.5). The arithmetic
action preserves each
rank-two source map, hence
their exact joint kernel
as well. The sphere factor
is the identity by
(D.46.10), whereas
D.45 computes a nonzero
degree-seven term for the
old action. This is an
explicit distinction under
the same marked quotient
map.

\emph{The entire relative-prime source and its parameter.}
For finitely many primes,
the actual image of \(N\)
in a product of finite
\(p\)-group quotients is
the full product, by the
coprime-order Goursat
argument. Passing to all
finite levels and primes
proves (D.46.2) with its
original source kernels.
Its factors have exactly
the common quotient \(Q\):
in any larger common
quotient retaining \(Q\),
the kernel would be both
a \(p\)-group and a
prime-to-\(p\) group.

Since \(\alpha_{h,2}\)
is the identity on \(Q\),
the product
\[
          \alpha^\dagger
              =(\alpha_{h,2},(1_{\mathcal P_p})_{p\ne2})
                                                               \tag{D.46.16}
\]
preserves the actual fibre
product and is a continuous
automorphism. It fixes \(t\).
Write the canonical normalized
Galois conjugators, in the
image of \(\widehat F\), as
\[
 \alpha_h(s)=\beta_hs\beta_h^{-1},\qquad
 \alpha_h(r)=A_hrA_h^{-1}.
\]
Put \(a_Q=q(A_h)\),
\(b_Q=q(\beta_h)\).
Both lie in \(q(F)\);
they centralize \(r_Q,s_Q\),
respectively, by (D.46.12).

For clarity, the coprime
fixed-element statement
used next applies to
these literal cosets.
Let a finite group of
order prime to \(p\)
act on a profinite
group with invariant
pro-\(p\) kernel.
Every fixed element of
the quotient lifts to
a fixed element whenever
that acting group is the
cyclic order-two or
order-three subgroup
used here. To prove it,
lift the quotient element;
its conjugate of the
acting cyclic subgroup
and the original subgroup
are complements in the
same pro-\(p\)-by-cyclic
extension. Coprime
complement conjugacy
corrects the chosen
lift by the kernel.
Finite centre induction
and compactness give the
profinite statement, as
in D.42. This is a
coset lifting result,
not an assumption that
the quotient element
already lies in the
kernel.

Apply it to conjugation
by \(r\) on
\(\mathcal P_2\to Q\).
There is an actual
\[
 c_2\in C_{\mathcal P_2}(r),
                  \qquad q(c_2)=a_Q^{-1}.                     \tag{D.46.17}
\]
Set \(A_2^\dagger=A_{h,2}c_2\).
It has \(Q\)-coordinate
one and gives exactly
the same conjugation
of \(r\) as before.
For every odd prime
apply the same result
to conjugation by \(s\):
choose
\[
 b_p\in C_{\mathcal P_p}(s),\qquad
                         q(b_p)=b_Q.                         \tag{D.46.18}
\]
One choice at each prime
is a single compatible
element of the full
product in (D.46.2).
The actual simultaneous
conjugators are
\[
 A^\dagger=(A_{h,2}c_2,(1)_{p\ne2}),\qquad
 \beta^\dagger=(\beta_{h,2},(b_p)_{p\ne2}).                    \tag{D.46.19}
\]
Their common \(Q\)
coordinates are \(1\)
and \(b_Q\), so both
are actual elements of
\(\mathcal U\). Their
images lie in the
indicated \(F\) source,
because their \(Q\)
coordinates do.

The normalized parameter
is recovered from the
automorphism of \(s\),
not from an arbitrary
choice of its conjugator:
\[
 f^\dagger=s\beta^\dagger s(\beta^\dagger)^{-1}
                         =(f_{h,2},(1)_{p\ne2}).              \tag{D.46.20}
\]
In particular the change
(D.46.17) changes neither
this parameter nor the
two-component automorphism.
For each prime \(p\),
the source quotient
\(\mathcal P_p\) maps
to the actual relative
pro-\(p\) completion of
the original \(F\):
the restriction of
any such map to \(N\)
has pro-\(p\) image
and kills \(R_p\).
Therefore \(\mathcal U\)
retains all
\(\widehat F^{\rm ab}
=\prod_p\mathbb Z_p^2\).
Equation (D.46.20) has
zero coordinate there.
Every actual \(F\)
preimage of it consequently
lies in \(\widehat F'\).
This proves the parameter
claim with its source
map, and also shows
that its image in
every pure \(F^{(p)}\)
is one.

The order-two and
order-three equations,
and the fixed \(t\),
give the two multiplier-one
hexagons by D.25's exact
braid-source comparison.
Equivalently
\(\alpha^\dagger(s)=s f^\dagger\)
and
\(\alpha^\dagger(r)=(f^\dagger)^{-1}r\).
The automorphism is
already bijective by
(D.46.16), so the
generation condition is
also satisfied. Finite
characteristic refinements
give the stated finite
GT shadows: their
pullbacks to \(B_3\)
are normal subgroups
of \(PB_3\), the
multiplier is one,
and the same actual
commutator parameter
projects compatibly.

\emph{The framing is an invariant of the new joint action.}
The distinction between
(D.46.13) and (D.46.3)
is checked using the
actual D.25 maps.
On \(G_2\) the new
automorphism is the
arithmetic \(h\) action.
The first equality in
(D.25.6) forces the
\(Q\)-coordinate of
every normalized
\(s\)-conjugator to
be \(b_Q\).
On \(G_3\) the new
automorphism is the
identity. The second
equality in (D.25.6)
forces the \(Q\)
coordinate of every
normalized \(r\)-conjugator
to be one.
Thus the new finite
shadow itself determines
the coordinates
\[
                 (A_Q^\dagger,\beta_Q^\dagger,V_Q)
                                  =(1,b_Q,1).                 \tag{D.46.21}
\]
They have not been
assigned as independent
matrix-centralizer choices.
The normalized congruence
matrix is \(V_Q=1\):
the action on \(Q\) is
the identity, so this
matrix is central in
the projective congruence
group; its lower-right-one
normalization and
multiplier one force
the scalar matrix \(I\).
The actual formulas
\(T=A^{-1}V\) and
\(W=\beta^{-1}V\),
with the unique global
sign fixing the rational
order-three point for
\(T_4\), now give
\(T_4=I\) and retain
the arithmetic
\(W_3=[i]\).

The finite certificate
also checks every
possible arithmetic
framing used in this
argument. The maximal
relative-two quotients
of \(Q\) over the
sphere and elliptic
bases have orders
\(24\) and \(48\).
Their source kernels
have orders \(24,12\),
trivial intersection,
and product of index
two. The first quotient
is the actual canonical
mod-four \(S_4\), checked
on all \(576^2\)
source pairs. Hence
\[
               Q=Q_F\times_{C_2}Q_E,\qquad
                       |Q_F|=24,\ |Q_E|=48.                  \tag{D.46.22}
\]
There are three possible
\(a_Q\) with arithmetic
\(T_4=-I\).
There are two possible
\(b_Q\) with
\(W_3=[i]\) and trivial
image in \(Q_F\).
The latter condition
holds here in the full
source: the preimage of
\(\langle s\rangle\)
in the sphere relative
completion is the free
pro-two product
\(C_2\amalg\mathbb Z_2\).
Its marked discrete
subgroup has index
three in \(B\), one
order-two vertex group
and one free generator;
the full induced pro-two
topology identifies the
completion. Herfort--Ribes
\cite[Theorem B$'$]{A_HerfortRibes1985}
gives \(C(s)=\langle s\rangle\)
there. Any centralizer
element in the whole
sphere source projects
into \(C_{S_3}(s)=\langle s\rangle\),
so lies in that subgroup.
As \(\beta_{h}\) belongs
to \(F\), the identity
sphere action forces
its projection to be
one. This proves the
stated condition on
\(b_Q\).

For all six resulting
possibilities the
certificate explicitly
multiplies \(a_Q\)
by \(a_Q^{-1}\),
retains \(b_Q\), and
computes the D.25
global-sign bits.
They change from
\((1,1)\) to
\((0,1)\) in every
case. This verifies
the precise framing
map without requiring
an unidentified choice
among the six cases.

\emph{The full-source centralizers and the remaining kernel.}
Herfort--Ribes
\cite[Theorem B$'$]{A_HerfortRibes1985}
also gives
\[
 C_{\widehat B}(r)=\langle r\rangle,\qquad
 C_{\widehat B}(s)=\langle s\rangle.                            \tag{D.46.23}
\]
Both cyclic images
meet \(q(F)\)
trivially. Every
nontrivial \(a_Q^{-1}\)
in (D.46.17) lies
in \(q(F)\), and
the two possible
\(b_Q\)'s in
(D.46.18) do too.
The program verifies
these containments
and exclusions for
all six cases.
Thus those relative
centralizer lifts
do not lift to
elements of the
corresponding full
centralizers.
This explains exactly
why the construction
does not alter the
same Galois automorphism
on the full source:
the odd relative
actions in
(D.46.16) have been
replaced by the
identity, and the
new parameter is
(D.46.20), not the
entire \(f_h\).
It does not prove
nonexistence of a
different full
automorphism inducing
this new parameter.

Taking the projection
kernels in (D.46.2)
gives (D.46.8).
For a concrete test
of the remaining
source kernel, the
program enumerates
all \(5040\)
permutations generated
by (D.46.9), with
odd \(s\) and even
\(r\).
The group \(Q\) is
solvable: its normal
two-subgroup has
quotient \(S_3\times C_3\).
Every finite image
of \(\mathcal U\)
is an extension of
a nilpotent group
by a quotient of
\(Q\), hence is
solvable.
The quotient of
\(S_7\) by the
image of
\(R_{\rm nil}\)
is therefore
solvable, so that
image contains
\(A_7\). Conversely
\(R_{\rm nil}\subset N
\subset B_{\rm even}\)
makes its image
even. The image
is exactly \(A_7\),
as asserted. This
is an actual
remaining full-source
kernel, despite the
entire relative-prime
lift proved above.
\end{proof}
